% !TEX program = xelatex
\documentclass[11pt]{ctexart}
%=====================================================================
% theme-violet.tex —— 讲义精讲 · 紫花主题
% 用法：\documentclass[11pt]{ctexart}
%       \input{assets/theme-violet}
%       \begin{document} ... \end{document}
% 只用 ctexart 自带能力 + xcolor/enumitem/fancyhdr/pifont/amsmath，
% 不依赖 tikz/tcolorbox，保证 MiKTeX 可直接编译。
%=====================================================================

\usepackage[a4paper,left=2.0cm,right=2.0cm,top=2.1cm,bottom=2.3cm,%
            headsep=12pt,footskip=22pt]{geometry}
\usepackage{amsmath,amssymb}
\usepackage{needspace}
\usepackage{xcolor}
\usepackage{array}
\usepackage{enumitem}
\usepackage{pifont}
\usepackage{fancyhdr}
\usepackage[hidelinks]{hyperref}

%---------------------------------------------------------------- 配色
\definecolor{violet}{RGB}{107,78,143}      % 主色：紫
\definecolor{violetdeep}{RGB}{63,44,92}    % 深紫：标题/强调
\definecolor{petal}{RGB}{183,159,208}      % 花瓣：装饰线、花饰
\definecolor{bandviolet}{RGB}{247,243,251} % 浅紫底：档案/提示框
\definecolor{bandcream}{RGB}{252,249,242}  % 米白底：易错/注意框
\definecolor{rulepurple}{RGB}{217,207,230} % 分隔线
\definecolor{inkgray}{RGB}{92,96,110}      % 次要文字

%---------------------------------------------------------------- 字体
% 有 Noto CJK 用 Noto，没有就退回 ctex 默认（Windows 下为宋体/黑体）
\IfFontExistsTF{Noto Serif CJK SC}{\setCJKmainfont{Noto Serif CJK SC}}{}
\IfFontExistsTF{Noto Sans CJK SC}{\setCJKsansfont{Noto Sans CJK SC}}{}

%---------------------------------------------------------------- 版式
\setlist{nosep,leftmargin=1.9em}
\linespread{1.28}
\setlength{\parskip}{3.5pt}
\setlength{\parindent}{0pt}
\setlength{\tabcolsep}{5pt}
\renewcommand{\arraystretch}{1.28}

%---------------------------------------------------------------- 花饰
% \ding{95} = 五瓣小花，\ding{96} = 星形花；都来自 pifont，必定存在
\newcommand{\orn}{\textcolor{petal}{\ding{95}}}          % 行内小标
\newcommand{\ornstar}{\textcolor{petal}{\ding{96}}}      % 星形花
\newcommand{\flowerrule}{%                                % 居中花饰分隔线
  \par\vspace{7pt}%
  \noindent\makebox[\linewidth]{\color{petal}%
    \rule{0.20\linewidth}{0.6pt}\hspace{8pt}{\small\ding{95}}\hspace{4pt}%
    \raisebox{1.1pt}{\tiny\ding{95}}\hspace{4pt}{\small\ding{95}}\hspace{8pt}%
    \rule{0.20\linewidth}{0.6pt}}%
  \par\vspace{7pt}}

%---------------------------------------------------------------- 标题与节头
\newcommand{\doctitle}[3]{%   #1 主标题  #2 副标题  #3 元信息（可用 \\ 分行）
  \begin{center}
    {\color{petal}\small \orn\enspace #2\enspace \ornstar}\\[7pt]
    {\LARGE\bfseries\color{violetdeep} #1}\\[10pt]
    {\color{petal}\rule{0.55\linewidth}{1pt}}\\[4pt]
    {\small\color{inkgray} #3}
  \end{center}\vspace{2pt}}

\newcommand{\hdr}[2]{%   #1 章节号（如 §0）  #2 标题
  \par\vspace{12pt}%
  \noindent\colorbox{violetdeep}{\color{white}\bfseries\small\;#1\;}\hspace{7pt}%
  {\color{violetdeep}\large\bfseries #2}\par\nobreak
  \vspace{3pt}{\color{petal}\hrule height 1.2pt}\vspace{6pt}}

\newcommand{\subhdr}[1]{%
  \par\vspace{6pt}%
  \noindent{\color{violet}\bfseries #1}\par\vspace{1pt}}

%---------------------------------------------------------------- 提示框
\newcommand{\callout}[1]{% 浅紫底：讲义档案、读法说明
  \par\medskip\noindent
  \fcolorbox{rulepurple}{bandviolet}{\parbox{\dimexpr\linewidth-2\fboxsep-2\fboxrule}{\vspace{2pt}#1\vspace{2pt}}}%
  \par\medskip}

\newcommand{\warnbox}[1]{% 米白底：易错、边界条件
  \par\medskip\noindent
  \fcolorbox{rulepurple}{bandcream}{\parbox{\dimexpr\linewidth-2\fboxsep-2\fboxrule}{\vspace{2pt}#1\vspace{2pt}}}%
  \par\medskip}

\newcommand{\gapitem}[2]{% #1 缺口类型（紫） #2 内容
  \par\vspace{2pt}\noindent
  \hangindent=1.4em\hangafter=1
  {\color{petal}\small\ding{96}}\enspace{\bfseries\color{violetdeep}#1}\enspace #2\par}

% 题干块（\qstem）：每道例题／作业的「详解」开头必须先复述题干与出处。
% 起因：2026-10-25 第五章讲义把七道作业的题面集中放在一页、详解另起一页，
% 读者翻到「第 4 题」时看不到条件，等于一道没有题干的题目。
% 规则：详解块内自带题干，读者不翻页就能读懂题在问什么。
\newcommand{\qstem}[1]{%
  \par\medskip\noindent
  \fcolorbox{violet}{bandviolet}{%
    \parbox{\dimexpr\linewidth-2\fboxsep-2\fboxrule}{%
      \vspace{2pt}%
      {\small\bfseries\color{violet}题干}\hspace{0.8em}#1%
      \vspace{2pt}}}%
  \par\medskip}

% v2: 保护题干与开头，长题可跨页，续页补上下文。
\newcommand{\qbreak}{\needspace{6\baselineskip}}
% 独立练习：ID、来源、题干、考点、答案标签（通常 ans:ID）。
\newcommand{\practice}[5]{%
  \qbreak\subhdr{练习 #1}\qstem{（#2）#3}%
  \textbf{考点}：#4\par
  详解见附录第~\pageref{#5}~页。\par\medskip}
% 附录开头：ID、来源、条件摘要；自动产生 ans:ID 标签。
\newcommand{\answerstart}[3]{%
  \qbreak\subhdr{#1\quad 详解}\label{ans:#1}\qstem{（#2）#3}}
\newcommand{\answercontinue}[2]{%
  \needspace{5\baselineskip}\subhdr{#1\quad 续解}%
  {\small\color{inkgray}本页接续：#2}\par\medskip}
% 延迟提供默认值，让旧讲义可在导言区保留自己的 \newcommand{\step}。
\AtBeginDocument{\providecommand{\step}[1]{\needspace{3\baselineskip}\par\medskip\textbf{#1}\par}}

% 前提旁注：只放在「需要理解的那一步」后面，给陈述 + 本题代入，不重证。
\newcommand{\aside}[1]{%
  \par\vspace{1pt}\noindent
  \hangindent=1.6em\hangafter=1
  {\small\color{inkgray}%
    \hspace*{0.6em}{\color{petal}$\triangleright$}\hspace{2pt}#1}%
  \par\vspace{2pt}}

%---------------------------------------------------------------- 表格
% 本机 MiKTeX 的 colortbl 与 array 版本不匹配（\rowcolor 会报错），
% 故表格统一用「上粗紫线 + 表头 + 中细线」的三线式，不依赖 colortbl。
\newcommand{\tabtop}{\par\vspace{3pt}{\color{violet}\hrule height 1.2pt}\vspace{1pt}}
\newcommand{\tabmid}{\par\vspace{1pt}{\color{petal}\hrule height 0.5pt}\vspace{1pt}}
\newcommand{\tabbot}{\par\vspace{3pt}{\color{rulepurple}\hrule height 0.5pt}\vspace{8pt}}
\newcommand{\thead}[1]{{\bfseries\color{violetdeep}#1}}

%---------------------------------------------------------------- 页眉页脚
\pagestyle{fancy}
\fancyhf{}
\renewcommand{\headrulewidth}{0.4pt}
\renewcommand{\headrule}{\hbox to\headwidth{\color{rulepurple}\leaders\hrule height \headrulewidth\hfill}}
\fancyhead[L]{\small\color{inkgray}\songti \doctag}
\fancyhead[R]{\small\color{inkgray}\songti \coursetag}
\fancyfoot[C]{\small\color{inkgray}\orn\enspace{\color{violet}\thepage}\enspace\orn}

% 每篇讲义在导言区设这两个值
\newcommand{\doctag}{}
\newcommand{\coursetag}{}
\newcommand{\settags}[2]{\renewcommand{\doctag}{#1}\renewcommand{\coursetag}{#2}}

%---------------------------------------------------------------- 数学记号
\newcommand{\vect}[1]{\boldsymbol{#1}}
\newcommand{\dd}{\,\mathrm{d}}
\newcommand{\ddx}{\frac{\mathrm{d}}{\mathrm{d}x}}
\addtolength{\jot}{2pt}



\setlength{\headheight}{22pt}
\usepackage{tikz}
\usetikzlibrary{positioning,arrows.meta,calc}
\settags{实验 5 · 上课执行单}{电工电子技术}
\let\orighdr\hdr
\renewcommand{\hdr}[2]{\needspace{4\baselineskip}\orighdr{#1}{#2}}
\let\origsubhdr\subhdr
\renewcommand{\subhdr}[1]{\needspace{3\baselineskip}\origsubhdr{#1}}

\begin{document}

\doctitle{实验 5 上课执行单：单管放大电路}
{2026-11-05（周四）15:30--17:45 · 综合一号教学楼 317 · 10 分 · 带预习报告和教材}
{这是现场用的排程卡，不是讲义。理论、六张表的完整条件与英文报告写法见《实验 5 单管放大电路 实验讲义》与《速查卡》\ \textperiodcentered\ 2026-11-05}

\callout{\textbf{这张纸怎么用。} §1 是出门清单，§2 是两小时十五分的排程（照时刻走，不要按“做完一项再做下一项”自由发挥），§3 是六张表的先后与条件，§4 是现场最容易翻的三处，§5 是时间不够时的取舍，§6 是收尾，§7 是当晚写报告的骨架。讲义里已有的公式这里不重复。}

%=============================================================
\hdr{§1}{出门前 5 分钟：带齐这五样}

\begin{center}
\small
\begin{tabular}{@{}p{0.6cm}p{4.6cm}p{9.4cm}@{}}
\hline
\thead{$\square$} & \thead{物件} & \thead{说明} \\
\hline
$\square$ & 预习报告（纸质） & 册上原文：\textbf{没有预习报告或预习报告不合格者不允许做实验}。按教材 §5.2.2 的四组问题写，哪怕写得粗也要有。 \\
$\square$ & 教材《电工学实验教程》第 2 版 & 册上原文：\textbf{不带教材不允许进实验室}。指到 §5.2（书 p.147--154）。 \\
$\square$ & 讲义与速查卡（可选） & 速查卡上有四组静点表、三个公式、报告六栏；讲义 20 页不必带进实验室。 \\
$\square$ & 计算器 + 笔 + 尺 & 现场要算 $I_B$、$I_C$、$U_{CE}$、$A_u$、$r_i$、$r_o$，并画波形。 \\
$\square$ & 手机（本题不禁） & \textbf{今天这场是普通实验，不是考试}；手机禁用在 2026-11-10 的实验 6。 \\
\hline
\end{tabular}
\end{center}

\warnbox{\textbf{报告双面打印、所有原始数据须指导教师确认批改，否则无效}；抄袭数据或编造数据记 0 分。今天现场记下的数就是最终数，别回宿舍“补”。}

%=============================================================
\hdr{§2}{时间盒：135 分钟照这张表走}

\begin{center}
\small
\begin{tabular}{@{}p{1.9cm}p{8.0cm}p{4.5cm}@{}}
\hline
\thead{时刻} & \thead{做什么} & \thead{产出/判据} \\
\hline
15:30--15:40 & 交预习报告、听老师讲注意事项、打开箱面认端子（a、b1、b2、d1、d2、c、$S_1$、$S_2$、$R_{P1}$、$R_{P2}$） & 能在箱面上指出每一个端子 \\
15:40--15:55 & \textbf{认管子}（二极管挡判 $b,c,e$）＋ \textbf{断电查线}（电阻挡三查：$U_{CC}$ 到地不短、基极到地约 $R_{B1}+R_{P1}$、集电极到地约 $R_C$） & 报告里能写出实测管脚排列 \\
15:55--16:15 & \textbf{表 6-1}：四组偏置的静态工作点（$R_{P1}=0$、$S_1$ 断开、\textbf{输入对地短接}） & 四组 $U_B,U_E,U_C$ $\Rightarrow$ $I_B,I_C,U_{CE}$ \\
16:15--16:40 & \textbf{表 6-2}：四组偏置 $\times$ $S_1$ 开/合（共 8 个波形 + $U_o$），并选定第 3 组 & 每格填“不失真/饱和削波/截止削波/双向” \\
\hline
\end{tabular}
\end{center}

\begin{center}
\small
\begin{tabular}{@{}p{1.9cm}p{8.0cm}p{4.5cm}@{}}
\hline
\thead{时刻} & \thead{做什么} & \thead{产出/判据} \\
\hline
16:40--16:55 & \textbf{表 6-3}：第 3 组、$S_1$ 闭合，$R_L=3/6/13\,\mathrm{k}\Omega$ 三档 & 三行 $A_u$ 与 $20\lg|A_u|$ \\
16:55--17:10 & \textbf{表 6-5} 输入电阻（串 $R$）＋ \textbf{表 6-6} 输出电阻（开路/带载） & 两次测量的 $U_i$ 必须相同 \\
17:10--17:30 & \textbf{表 6-4}：频率响应 16 点（先低频段，再中频，再高频） & 中频 $U_o$ 与 0.707 判据 \\
17:30--17:40 & 请老师在原始数据上签字/批改；抄走全部数据；仪器归零、关机、拆线 & 数据带走、仪器断电 \\
17:40--17:45 & 收拾台面、交现场可交的东西 & 离开 \\
\hline
\end{tabular}
\end{center}

\callout{\textbf{为什么把表 6-4 放最后}：它要改 16 次频率、每次等读数稳定，最费时间；而它是唯一可能“测不出上限”的一张（表只到 $20\,\mathrm{kHz}$）。万一时间被前面吃掉，砍它比砍表 6-1/6-2/6-3 划算——后三张是讨论题的依据。}

%=============================================================
\hdr{§3}{六张表的先后与“先写理论值”}

\noindent 每张表动笔前先在草稿上写好条件行（$U_{CC}$、$R_{P1}$、$S_1$、$S_2$、输入幅度与频率），再填数；条件不写，数据没法复算，老师也会问。

\begin{center}
\small
\begin{tabular}{@{}p{1.3cm}p{4.4cm}p{2.4cm}p{6.6cm}@{}}
\hline
\thead{表} & \thead{条件} & \thead{行数} & \thead{先写什么理论值（讲义 §1.3、§5.5）} \\
\hline
6-1 & $U_{CC}=12\,\mathrm{V}$、$R_{P1}=0$、$S_1$ 断开、输入短路 & 4 组 & 第 1 组接近饱和（$U_{CE}\approx0.3\,\mathrm{V}$）；第 3 组 $I_C\approx1.17\,\mathrm{mA}$、$U_{CE}\approx7.10\,\mathrm{V}$ \\
6-2 & 输入 $10\,\mathrm{mV}$、$1\,\mathrm{kHz}$、$S_2$ 断开 & 4 组 $\times$ 2 & 每组的“预期状态”（第 1 组必削波、第 4 组近截止） \\
6-3 & 第 3 组、$S_1$ 闭合、$10\,\mathrm{mV}$、$1\,\mathrm{kHz}$ & 3 档 & $R_L$ 越小 $|A_u|$ 越小：三档单调下降 \\
6-4 & 固定 $U_i$，只改 $f$ & 16 点 & 中频段 $|A_u|$ 的最大值；$f_H$ 可能测不到 \\
6-5 & 串 $R$（$1\sim2\,\mathrm{k}\Omega$）、$S_1$ 闭合与断开 & 2 行 & $S_1$ 断开（有反馈）$r_i$ 更大 \\
6-6 & $R_L$ 已知、$S_1$ 闭合与断开 & 2 行 & $r_o$ 约 $R_C$ 量级（几 $\mathrm{k}\Omega$） \\
\hline
\end{tabular}
\end{center}

\aside{表 6-1 的 $I_B$ 那一栏，教材表 5.2.2 把单位印成了 “V”，应为 $\mu$A。照错的表头抄进报告，等于告诉老师你没核对。}

%=============================================================
\hdr{§4}{现场最容易翻的三处}

\warnbox{\textbf{一、$S_1$ 的方向记反。}$S_1$ 与 $R_{E2}$ \textbf{并联}：\textbf{闭合 = 短接 $R_{E2}$ = 没有这级负反馈}。如果记成“闭合才有反馈”，表 6-2 的八格解读会整体反过来，讨论题也跟着错。记不住的记法：$S_1$ 闭合 $\Rightarrow$ 发射极只剩 $R_{E1}$ $\Rightarrow$ $|A_u|$ 变大。}

\warnbox{\textbf{二、静态测量前没把输入对地短接、$R_{P1}$ 没归零。}教材 p.148 原话是“应先将放大电路的输入端与地短接”。输入悬空时示波器/毫伏表的拾取会让 $U_B$ 抖动；$R_{P1}$ 不在零位，四组偏置就不在同一条件下，比较失去意义。}

\warnbox{\textbf{三、表 6-5/6-6 的两次测量换了输入幅度。}$r_i=\dfrac{U_i}{U_s-U_i}R$ 与 $r_o=\left(\dfrac{U_{oc}}{U_{oL}}-1\right)R_L$ 都默认“这两次是同一个源”。改 $R$ 或合上 $S_2$ 之后，先把信号源幅度重新调到原来的 $U_i$，再读数；两次都不许失真。}

\callout{\textbf{还有两条现场纪律}：（一）示波器与毫伏表的地夹一律接端子 c（公共地），不要夹到电路中间某点；（二）$C_1$、$C_2$、$C_E$ 是电解电容，$+$ 端朝高电位（$C_1$ 朝基极、$C_2$ 朝集电极、$C_E$ 朝发射极），接反了静态点就会偏。}

%=============================================================
\hdr{§5}{时间不够时的取舍顺序}

\begin{enumerate}
  \item \textbf{必保}：表 6-1、表 6-2、表 6-3 —— 讨论题第 1--5 问全靠它们。
  \item \textbf{能保则保}：表 6-5、6-6 —— 每张只有两行，各 3 分钟；$r_i$ 也可用“电位器调到 $U_i=0.5U_s$、再量 $R_P$”的变体，更快。
  \item \textbf{可压缩}：表 6-4 只测低频 4 点（$10,50,100,200\,\mathrm{Hz}$）＋中频 $1\,\mathrm{kHz}$ ＋高频 3 点（$10,20\,\mathrm{kHz}$），报告里写明“受课时限制只测 8 点”。
  \item \textbf{不许做}：为了凑出 $f_H$ 编一个数。讲义 §5.6 已说明单管共射的上限通常在几百 $\mathrm{kHz}$，表只到 $20\,\mathrm{kHz}$，测不到是正常结论，如实写就得分。
\end{enumerate}

%=============================================================
\hdr{§6}{收尾：老师会看的两件事}

\begin{enumerate}
  \item \textbf{数据签字}：册上规则是“所有记录须指导教师确认批改，否则无效”。走之前把六张表拿给老师看一遍，请他在原始数据上签字或打勾。
  \item \textbf{关仪器}：老师在课上反复强调“记得实验完关所有仪器”。顺序：先关信号源输出，再关电源，再关示波器与毫伏表；$U_{CC}$ 断开后拆线；电解电容放电（用电阻或等一会儿）。
\end{enumerate}

%=============================================================
\hdr{§7}{当晚 40 分钟：报告骨架}

\begin{center}
\small
\begin{tabular}{@{}p{1.2cm}p{3.6cm}p{9.8cm}@{}}
\hline
\thead{分钟} & \thead{栏} & \thead{写什么} \\
\hline
0--5 & I. Objectives & 三句：测静态工作点；测 $A_u$、$r_i$、$r_o$、通频带；看负反馈（$S_1$）与负载的影响。 \\
5--15 & II. Principles & 分压偏置的静态公式、$r_{be}$ 与小信号模型、$A_u$、$r_i$、$r_o$、$f_{BW}$ 五组式子，各一行。 \\
15--25 & III. Data & 六张表照抄，\textbf{条件行一起抄}；表 6-2 每格补“状态”与波形草图。 \\
25--33 & IV. Discussion & 七问，每问 2--4 句，先结论后理由；第 4 问写“带载时先碰截止侧，$U_{op\text{-}p}\approx3.5\,\mathrm{V}$”。 \\
33--37 & V. Error & 电阻容差、$\beta$ 分散、毫伏表精度、$R_{P1}$ 读数、探头分压。 \\
37--40 & VI. Conclusion & 一句“实测与理论在 $X\%$ 内一致，$f_H$ 未测出”，不写长篇。 \\
\hline
\end{tabular}
\end{center}

\callout{\textbf{英文栏名与术语表}在讲义 §7.2、§7.4；\textbf{七问的参考答法}在讲义附录 B。报告要英文写，但数据表里的符号与单位照符号体系，不译。}

%=============================================================
\hdr{§8}{仍待现场确认的两件事}

\begin{enumerate}
  \item \textbf{箱面上 $R_{P1}$ 的接法与 $C_E$ 的跨接位置}：教材图 5.2.4 是原理图，箱面用跳线实现，可能不完全一致。通电前用电阻挡量 $U_{CC}$ 端到基极的阻值（应约为 $R_{B1}+R_{P1}$）、量 $C_E$ 两端是否只跨 $1\,\mathrm{k}\Omega$。
  \item \textbf{老师这次的检查方式}：预习报告是收走还是当面问、数据是逐个看还是走一圈。到场先观察一次，再决定第 15 分钟先做哪张表。
\end{enumerate}

\flowerrule

\noindent{\small 一句话总结：\textbf{先保表 6-1、6-2、6-3，再补 6-5、6-6，最后做 6-4；$S_1$ 闭合＝没有反馈；两次测量不许换输入幅度。}}

\clearpage

\noindent{\large\bfseries 附页：图 5.2.4 与端子（现场照这张接线）}

\vspace{4pt}

\begin{center}
\includegraphics[width=0.86\linewidth]{图-实验5-教材图524-无题注.png}

{\small 图 5.2.4\quad 共射极单管放大电路（\textbf{裁剪}自教材书 p.151）。a、b1、b2、d1、d2、c 是实验箱面上的接线柱，不是元件；$S_1$、$S_2$ 是箱面开关。}
\end{center}

\begin{center}
\footnotesize
\begin{tabular}{@{}p{1.5cm}p{6.4cm}p{6.8cm}@{}}
\hline
\thead{端子} & \thead{图上位置} & \thead{接错会怎样} \\
\hline
a & $R_{P1}$ 下端，跳线到 b1 或 b2 & 悬空则基极无偏置，$I_C\approx0$ \\
b1 / b2 & 上偏置 $30\,\mathrm{k}\Omega$ / $150\,\mathrm{k}\Omega$ 的上端 & 选错档，四组偏置就与表 6-1 对不上 \\
d1 / d2 & 下偏置 $39\,\mathrm{k}\Omega$ / $12\,\mathrm{k}\Omega$ 的接地端 & 悬空则基极被抬到接近 $U_{CC}$，饱和 \\
c & 公共地 & 多点接地会引入干扰，读数抖动 \\
$S_1$ & 与 $R_{E2}=200\,\Omega$ 并联 & \textbf{闭合＝短接 $R_{E2}$＝无负反馈} \\
$S_2$ & 串在 $R_{L1}=3\,\mathrm{k}\Omega$ 负载支路 & 断开测的是 $U_{oc}$，别当带载电压 \\
\hline
\end{tabular}
\end{center}

\warnbox{\textbf{报告里的电路图}：直接照这一页画（或把这一页打印出来贴）。英文引用句与手绘顺序在《实验 5 电路图页》，共两页；这张附页只给接线，不给写法。}

\end{document}
